% ==================================================================
% id:       cond-014
% topic:    הסתברות מותנית ואי-תלות
% subtopic: נוסחת ההסתברות השלמה ונוסחת בייס
% solution: full
% ==================================================================
בכד $A$ יש 8 כדורים לבנים, 7 אדומים ו-5 שחורים. בכד $B$ יש 4 לבנים, 3 אדומים ו-2 שחורים. מוציאים באקראי כדור מכד $A$ ומעבירים אותו לכד $B$. אחר כך מוציאים באקראי כדור מכד $B$. מה ההסתברות שהכדור שיצא מכד $B$ אדום?
\hint{חלקו לשני מקרים, לפי הצבע של הכדור שהועבר: אדום או לא אדום.}
\begin{solbox}
נסמן $R_1$ = ״הכדור שהועבר מכד $A$ אדום״, ו-$R_2$ = ״הכדור שהוצא מכד $B$ אדום״. המאורעות $R_1$ ו-$R_1^C$ הם חלוקה של $\Omega$, ולכן לפי נוסחת ההסתברות השלמה
\[
\PP(R_2)=\PP(R_2\mid R_1)\,\PP(R_1)+\PP(R_2\mid R_1^C)\,\PP(R_1^C).
\]
נחשב כל אחד מהגורמים:
\begin{itemize}
\item בכד $A$ יש $8+7+5=20$ כדורים, מהם 7 אדומים, ולכן $\PP(R_1)=\frac7{20}$ ו-$\PP(R_1^C)=\frac{13}{20}$.
\item אחרי ההעברה יש בכד $B$ $9+1=10$ כדורים. אם הכדור שהועבר אדום, יש בכד $B$ $3+1=4$ אדומים, ולכן $\PP(R_2\mid R_1)=\frac4{10}$.
\item אם הכדור שהועבר לא אדום, נשארים בכד $B$ 3 אדומים, ולכן $\PP(R_2\mid R_1^C)=\frac3{10}$.
\end{itemize}
נציב:
\[
\PP(R_2)=\frac4{10}\cdot\frac7{20}+\frac3{10}\cdot\frac{13}{20}=\frac{28+39}{200}=\frac{67}{200}=\ans{0.335}.
\]
\end{solbox}
